% Part: normal-modal-logic
% Chapter: axioms-systems
% Section: derived-rules

\documentclass[../../../include/open-logic-section]{subfiles}

\begin{document}

\olfileid[id]{nml}{prf}{der}

\olsection{Aturan Turunan}

Mencari dan menulis !!{derivation}s jelas sulit, merepotkan, dan berulang.
Sebagai contoh, sangat sering kita ingin beralih dari $!A\lif !B$ ke
$\Box !A\lif\Box !B$, yakni menerapkan aturan~\RK. Hal itu memerlukan
penerapan~\Nec, lalu pencatatan instansiasi~\Ax{K} yang tepat, kemudian
penerapan~\MP. Beralih dari $!A\lif !B$ dan $!B\lif !C$ ke $!A\lif !C$
memerlukan pencatatan instansiasi tautologis (yang panjang)
\[
(!A \lif !B) \lif ((!B \lif !C) \lif (!A \lif !C))
\]
dan penerapan~\MP{} dua kali. Sering kali kita ingin mengganti suatu
sub-!!{formula} dengan formula yang kita ketahui ekuivalen dengannya, misalnya
\iftag{prvDiamond}{$\Diamond !A$ dengan $\lnot\Box\lnot !A$}{$\Box !A$
dengan $\lnot\Diamond\lnot !A$}, atau $\lnot\lnot !A$ dengan~$!A$. Jadi,
alih-alih menuliskan !!{derivation} yang sebenarnya, lebih mudah mencatat
mengapa langkah-langkah antara itu !!{derivable}. Untuk tujuan tersebut,
marilah kita kumpulkan beberapa fakta tentang !!{derivability}.

\begin{prop}
  Jika $\Log{K}\Proves !A_1$, \dots, $\Log{K}\Proves !A_n$, dan $!B$
  mengikuti dari $!A_1$, \dots, $!A_n$ melalui logika proposisional, maka
  $\Log{K}\Proves !B$.
\end{prop}

\begin{proof}
  Jika $!B$ mengikuti dari $!A_1$, \dots,~$!A_n$ melalui logika
  proposisional, maka
  \[
  !A_1 \lif (!A_2 \lif \cdots (!A_n \lif !B)\dots)
  \]
  merupakan instansiasi tautologis. Menerapkan~\MP{} sebanyak $n$ kali
  menghasilkan !!a{derivation} dari~$!B$.
\end{proof}

Kita akan menandai penggunaan proposisi ini dengan~\PL.

\begin{prop}
  Jika $\Log{K}\Proves !A_1\lif(!A_2\lif\cdots(!A_{n-1}\lif
  !A_n)\dots)$, maka $\Log{K}\Proves\Box !A_1\lif(\Box !A_2\lif
  \cdots(\Box !A_{n-1}\lif\Box !A_n)\dots)$.
\end{prop}

\begin{proof}
  Melalui induksi pada~$n$, persis seperti dalam pembuktian
  \olref[nor]{prop:rk}.
\end{proof}

Kita akan menandai penggunaan proposisi ini dengan~\RK. Mari kita gambarkan
bagaimana hasil-hasil ini mempermudah pembuktian hasil !!{derivability}.

\begin{prop}
  $\Log{K}\Proves(\Box!A\land\Box!B)\lif\Box(!A\land !B)$
\end{prop}

\begin{proof}
  \begin{derivation}
    1. & $\Log{K} \Proves !A \lif (!B \lif (!A \land !B))$ & \Taut \\
    2. & $\Log{K} \Proves \Box!A \lif (\Box !B \lif \Box(!A \land !B)))$ & \RK, 1\\
    3. & $\Log{K} \Proves (\Box!A \land \Box!B) \lif \Box(!A \land !B)$ & \PL, 2
  \end{derivation}
\end{proof}

\begin{prop}\ollabel{prop:rewriting}
  Jika $\Log{K}\Proves !A\liff !B$ dan
  $\Log{K}\Proves\Subst{!C}{!A}{q}$, maka
  $\Log{K}\Proves\Subst{!C}{!B}{q}$.
\end{prop}

\begin{proof}
  Latihan.
\end{proof}

\begin{prob}
  Buktikan \olref[nml][prf][der]{prop:rewriting} dengan membuktikan, melalui
  induksi pada kompleksitas~$!C$, bahwa jika $\Log{K}\Proves !A\liff !B$,
  maka $\Log{K}\Proves\Subst{!C}{!A}{q}\liff\Subst{!C}{!B}{q}$.
\end{prob}

Proposisi ini terutama berguna ketika kita ingin mengubah $\Diamond$ menjadi
$\Box$ (atau sebaliknya), atau menghapus negasi ganda di dalam !!a{formula}.
Dalam pembahasan berikut, kita akan menandai penerapan
\olref{prop:rewriting} dengan ``$!B$ untuk~$!A$'' ketika kita menulis ulang
!!a{formula}~$!C(!A)$ menjadi~$!C(!B)$. Dengan kata lain, ``$!B$ untuk~$!A$''
menyingkat:
  \begin{derivation}
    & $\Proves !C(!A)$ \\
    & $\Proves !A \liff !B$\\
    & $\Proves !C(!B)$ & menurut \olref{prop:rewriting}
  \end{derivation}
Sebagai contoh:

\begin{prop}
  $\Log{K}\Proves\lnot\Box p\lif\Diamond\lnot p$
\end{prop}

\begin{proof}
  \iftag{prvDiamond}{% Diamond is primitive
  \begin{derivation}
    1. & $\Log{K} \Proves \Diamond \lnot p \liff \lnot\Box\lnot\lnot p$ &
    \Dual\\
    2. & $\Log{K} \Proves \lnot \Box \lnot\lnot p \lif \Diamond \lnot p$ & \PL, 1\\
    3. & $\Log{K} \Proves \lnot\Box p \lif \Diamond \lnot p$ & $p$ untuk $\lnot\lnot p$\\
  \end{derivation}
  }{% Diamond is defined
  \begin{derivation}
    1. & $\Log{K} \Proves \lnot \Box \lnot\lnot p \lif \Diamond \lnot p$ & \Taut\\
    2. & $\Log{K} \Proves \lnot\Box p \lif \Diamond \lnot p$ & $p$ untuk $\lnot\lnot p$\\
  \end{derivation}
  Formula pada baris~$1$ merupakan instansiasi $\lnot q\lif\lnot q$, dengan
  $\Box\lnot\lnot p$ disubstitusikan untuk~$q$. $\Diamond\lnot p$ dan
  $\lnot\Box\lnot\lnot p$ identik, sebab $\Diamond$ didefinisikan sebagai
  $\lnot\Box\lnot$.}
\end{proof}

Dalam !!{derivation} di atas, langkah terakhir ``$p$ untuk $\lnot\lnot p$''
merupakan singkatan bagi
  \begin{derivation}
    & $\Log{K} \Proves \lnot \Box \lnot\lnot p \lif \Diamond \lnot p$\\
    & $\Log{K} \Proves \lnot\lnot p \liff p$ & \Taut\\
    & $\Log{K} \Proves \lnot\Box p \lif \Diamond \lnot p$ & menurut \olref{prop:rewriting}
  \end{derivation}
Peran $!C(q)$, $!A$, dan~$!B$ dalam \olref{prop:rewriting} di sini dimainkan,
berturut-turut, oleh $\lnot\Box q\lif\Diamond\lnot p$,
$\lnot\lnot p$, dan~$p$.

Jika !!a{formula} memuat sub-!!{formula}~$\lnot\Diamond !A$, kita dapat
menggantinya dengan $\Box\lnot !A$ menggunakan \olref{prop:rewriting}, sebab
$\Log{K}\Proves\lnot\Diamond !A\liff\Box\lnot !A$. Kita akan menandai
penggantian ini dan penggantian serupa cukup dengan ``$\Box\lnot$ untuk
$\lnot\Diamond$.''

Proposisi berikut membenarkan bahwa kita dapat menetapkan hasil
!!{derivability} secara skematis. Sebagai contoh, proposisi sebelumnya tidak
hanya menetapkan bahwa $\Log{K}\Proves\lnot\Box p\lif\Diamond\lnot p$, tetapi
juga $\Log{K}\Proves\lnot\Box !A\lif\Diamond\lnot !A$ bagi sebarang~$!A$.

\begin{prop}
  Jika $!A$ merupakan instansiasi substitusi dari~$!B$ dan
  $\Log{K}\Proves !B$, maka $\Log{K}\Proves !A$.
\end{prop}

\begin{proof}
  Membosankan tetapi rutin untuk memeriksa (melalui induksi pada panjang
  !!{derivation} dari~$!B$) bahwa penerapan suatu substitusi pada seluruh
  !!{derivation} juga menghasilkan !!a{derivation} yang benar. Secara khusus,
  instansiasi substitusi dari instansiasi tautologis juga merupakan instansiasi
  tautologis; instansiasi substitusi dari instansiasi
  \iftag{prvDiamond}{\Dual{} dan~}{}\Ax{K} juga merupakan instansiasi
  \iftag{prvDiamond}{\Dual{} dan~}{}\Ax{K}; dan penerapan~\MP{} dan~\Nec{}
  tetap benar ketika !!{formula}s disubstitusikan bagi
  !!{propositional variable}s dalam premis maupun kesimpulannya.
\end{proof}

\end{document}
